\documentclass[11pt]{article}
\usepackage[margin=1in]{geometry}
\usepackage{amsmath,amssymb,amsthm,booktabs}
\usepackage[hidelinks]{hyperref}
\usepackage{xurl}
\setlength{\emergencystretch}{2em}
\newtheorem{theorem}{Theorem}
\newtheorem{lemma}[theorem]{Lemma}
\newtheorem{corollary}[theorem]{Corollary}
\newcommand{\mr}{\operatorname{mr}}
\newcommand{\br}{\operatorname{br}}
\newcommand{\bi}{\operatorname{bi}}
\title{Prime braid families with sharp alternating-group meridional bounds}
\author{Written proofs and exact literal certificates}
\date{17 November 2026}
\begin{document}\maketitle

For every integer $b\ge5$ we give infinitely many distinct prime knots
with meridional rank, bridge number and braid index all equal to $b$.
Their groups map onto $A_{2b}$, with meridians mapping to double
transpositions. This target has ordinary generating rank two, whereas
its specified meridian class has generating rank $b$. The construction
uses explicit braid words, a uniform group-generation argument and the
visual primeness theorem for homogeneous braids. The genus and minimal
crossing number are computed exactly.

\section{The knots and the conclusions}

Put
\begin{equation}\label{eq:beta}
 \beta=\sigma_1\sigma_2\sigma_3^{-2}\sigma_2\sigma_3^{-2}
        \sigma_2\sigma_1^2\sigma_2^3\sigma_3^{-1}\sigma_2^3\in B_4,
\end{equation}
and, for $b\ge5$, put
\[
 W_b=\beta\,\sigma_4^3\sigma_5^3\cdots\sigma_{b-1}^3\in B_b,
 \qquad K_{b,n}=\widehat{W_b^n}.
\]
The braid convention is fixed by the signed literal PD certificate in
Section~\ref{sec:literal}. Reversing the crossing convention gives the
mirror, with the same stated invariants. Write $\br(K)$ for bridge
number, $\bi(K)$ for braid index, and $c(K)$ for minimal crossing number.

\begin{theorem}\label{thm:family}
Let $b\ge5$, $n\ge2$ and $\gcd(n,b)=1$. Then $K_{b,n}$ is a prime knot and
\[
 \mr(K_{b,n})=\br(K_{b,n})=\bi(K_{b,n})=b,
\]
\[
 g(K_{b,n})=\frac{n(3b+5)-b+1}{2},\qquad
 c(K_{b,n})=n(3b+5).
\]
There is an epimorphism
$\pi_1(S^3\setminus K_{b,n})\twoheadrightarrow A_{2b}$ taking every
meridian to a double transposition. For each fixed $b$ these formulas
give infinitely many distinct prime knots.
\end{theorem}

This theorem makes no assertion about the augmentation rank of these
powers, or about absence of other kinds of quotients. It does not require
a uniform hyperbolicity claim. The argument below also proves the
four-strand variant and a separate connected-sum statement.

\section{Ordinary rank and meridian-class rank}

Permutations act on $\{0,\ldots,N-1\}$ and $pq$ means $p\circ q$.
For a conjugacy class $C$ in a finite group $G$, its class generating rank
is the least number of elements of $C$ that generate $G$.

\begin{lemma}\label{lem:ordinary}
For even $N\ge4$ the two permutations
\[
 a_N=(0\,1\,\cdots\,N-2),\qquad b_N=(0\,1\,N-1)
\]
generate $A_N$. Thus its ordinary generating rank is two.
\end{lemma}
\begin{proof}
Both permutations are even. The conjugates of $b_N$ by powers of $a_N$
are the three-cycles $(k\,k+1\,N-1)$, with the first two entries read
cyclically among $0,\ldots,N-2$.
For any connected graph on vertices other than a fixed vertex $u$, the
three-cycles $(u\,i\,j)$ associated with its edges generate $A_N$.
Indeed, conjugating $(u\,y\,z)$ by $(u\,x\,y)$ gives $(u\,z\,x)$.
Edge completion along paths therefore gives all $(u\,i\,j)$. Moreover,
\[
 (u\,a\,b)(u\,b\,c)=(a\,b\,c),
\]
so the subgroup contains every three-cycle. Every even permutation is
a product of three-cycles. This proves generation; the nonabelian group
$A_N$ is not cyclic.
\end{proof}

Consider the following double transpositions on $2b$ points:
\begin{equation}\label{eq:generators}
\begin{aligned}
 x_0&=(0\,1)(2\,3),& x_1&=(1\,4)(2\,3),\\
 x_2&=(0\,5)(1\,6),& x_3&=(1\,2)(6\,7),\\
 x_j&=(1\,2j)(6\,2j+1)&& (4\le j\le b-1).
\end{aligned}
\end{equation}

\begin{theorem}\label{thm:targets}
For every $b\ge4$ the tuple $(x_0,\ldots,x_{b-1})$ generates $A_{2b}$.
The double transpositions form one $A_{2b}$ conjugacy class, of size
$3\binom{2b}{4}$ and class generating rank exactly $b$.
\end{theorem}
\begin{proof}
The base case has short directly checkable identities:
\begin{align*}
 x_3x_0x_1x_3x_0x_2x_0x_1x_2&=(0\,1\,2\,3\,4\,5\,6),\\
 x_0x_3x_1x_2x_3x_1x_3x_2x_1x_3&=(0\,1\,7).
\end{align*}
Lemma~\ref{lem:ordinary} gives $A_8$.

Here is the induction without enumerating any large group. Suppose the
subgroup already contains $A_m$ on the old points, where $m\ge8$, and
adjoin $z=(a\,u)(c\,v)$, with $a,c$ distinct old points and $u,v$ new.
Conjugation by an even permutation of the old points fixing $c$ and
sending $a$ to a different old point $d$ gives
$z'=(d\,u)(c\,v)$. Then
\[
                         zz'=(a\,u\,d).
\]
Conjugating by $A_m$, which is transitive on ordered pairs of distinct
old points, gives all three-cycles containing $u$ and two old points.
Together with $A_m$ these generate $A_{m+1}$ on the old points and $u$.
Now conjugate $z$ by an even permutation fixing $a,u$ and sending $c$
to a different old point $e$. This gives $z''=(a\,u)(e\,v)$, and
\[
                         zz''=(c\,v\,e).
\]
The same argument gives $A_{m+2}$. All generators are even, so the group
is exactly $A_{m+2}$. Taking $a=1,c=6,u=2j,v=2j+1$ proves the induction
for~\eqref{eq:generators}. The required conjugators can, for example, be
chosen as $(1\,0\,2)$ and $(6\,0\,2)$ respectively.

There are $3\binom{2b}{4}$ double transpositions in $S_{2b}$. They are
conjugate in $A_{2b}$ too: multiplying any odd conjugator by one
transposition factor centralizing its source makes the conjugator even.
For a family of $r$ double transpositions, draw their two transposition
edges per generator. Each generator preserves every component of this
graph. A transitive subgroup on $2b$ points therefore requires a
connected graph, so $2r\ge2b-1$ and $r\ge b$. The displayed generating
tuple realizes equality.
\end{proof}

The distinction between ordinary rank and class generating rank is a
property of the finite target. Ordinary rank two of $A_{2b}$ is not a
claim about the ordinary rank of the knot group.

\section{The quotient and the sharp bridge bound}

For involutions $u,v$, the braid coloring convention is
\begin{equation}\label{eq:hurwitz}
 \sigma_i:(u,v)\longmapsto(v,vuv),\qquad
 \sigma_i^{-1}:(u,v)\longmapsto(uvu,u)
\end{equation}
on the two affected strands. These are inverse operations. They are
exactly the two Wirtinger crossing rules for involutory images.

\begin{lemma}\label{lem:closure}
For every $b\ge4$, the braid $W_b$ returns the tuple
$(x_0,\ldots,x_{b-1})$ to itself, with $W_4=\beta$. Its strand
permutation is a $b$-cycle.
\end{lemma}
\begin{proof}
Applying the seventeen signed letters of~\eqref{eq:beta} using
\eqref{eq:hurwitz} returns the first four entries exactly. This is a
finite identity of permutations on eight points; the supplement checks
all seventeen steps, independently of topology software.
For $j\ge4$, the product $x_{j-1}x_j$ is a product of two disjoint
three-cycles. This is immediate for $x_3=(1\,2)(6\,7)$ and
$x_4=(1\,8)(6\,9)$, and for the later pairs from their common factors
at the points $1$ and $6$. Thus $(x_{j-1}x_j)^3=1$.
For involutions $u,v$ with $(uv)^3=1$, the braid relation $uvu=vuv$
shows directly that three applications of~\eqref{eq:hurwitz} return
$(u,v)$. Consequently every appended $\sigma_j^3$ returns its affected
pair, so $W_b$ returns the entire tuple.

Ignoring signs, the seventeen letters of $\beta$ give a four-cycle on
the strands. Appending $\sigma_j^3$ composes with the adjacent
transposition inserting the new strand $j+1$ in that cycle. Induction
gives a $b$-cycle. Its $n$th power has $\gcd(n,b)$ cycles, which count
the components of the braid closure.
\end{proof}

\begin{proof}[Meridional and bridge assertions of Theorem~\ref{thm:family}]
The returned tuple supplies consistent arc images for the closure of
$W_b^n$. A braid crossing enforces a Wirtinger relation, and the return
condition enforces every closing relation. Hence it defines a
homomorphism of the actual knot group. Its image contains all $x_j$,
so Theorem~\ref{thm:targets} proves surjectivity onto $A_{2b}$.
All meridians map into the double-transposition class, including after
inversion, since its elements are involutions. A meridional generating
family must therefore have at least $b$ elements by that theorem.

A closed $b$-braid has a bridge presentation with at most $b$ bridges:
place one maximum on each closing arc and make the braid strands
monotone between them and the corresponding minima. The elementary
bridge presentation of the group is generated by these bridge
meridians. Thus
\[
 b\le\mr(K_{b,n})\le\br(K_{b,n})\le\bi(K_{b,n})\le b.
\]
This proves all three equalities without inferring bridge number merely
from a meridional generating set.
\end{proof}

\section{Primeness, genus and minimal crossings}

A braid word is homogeneous if each generator occurs with only one
sign. In our words $\sigma_1,\sigma_2$ are positive, $\sigma_3$ is
negative, and every later generator is positive.
For $2\le i\le b-1$, project a word to the letters
$\sigma_{i-1}^{\pm1},\sigma_i^{\pm1}$. Its seesaw number $g_i$ is the
number of maximal nonempty consecutive blocks with the same generator
index. This is the \emph{linear} word definition in
\cite[Section~5.2, printed p.~20]{flo}; the last and first blocks are not
identified.

\begin{lemma}\label{lem:seesaws}
For $b\ge5$, $n\ge1$, the seesaw numbers of $W_b^n$ are
\[
 g_2=4n,\qquad g_3=6n+1,\qquad g_i=2n\quad(4\le i\le b-1).
\]
For $b=4$ the first two formulas hold.
\end{lemma}
\begin{proof}
The projection of one copy to the indices $1,2$ is
$\sigma_1\sigma_2^3\sigma_1^2\sigma_2^6$. Its four runs do not merge
across copies. The projection to the indices $2,3$ is
\[
 \sigma_2\sigma_3^{-2}\sigma_2\sigma_3^{-2}
 \sigma_2^4\sigma_3^{-1}\sigma_2^3.
\]
It has seven runs and both ends have index $2$, so exactly one run is
lost at each of the $n-1$ joins. The indices $3,4$ give
$\sigma_3^{-5}\sigma_4^3$; every later adjacent pair gives
$\sigma_{i-1}^3\sigma_i^3$. These have two runs, with no merges across
copies.
\end{proof}

\begin{proof}[Primeness assertion of Theorem~\ref{thm:family}]
The closure is a knot by Lemma~\ref{lem:closure}. Feller--Lewark--Orbegozo
Rodriguez prove that if the closure of a homogeneous braid is not prime,
its displayed diagram has a decomposition circle giving a decomposition
sphere \cite[Theorem~1.1, printed p.~1]{flo}.
Their Lemma~5.1(i), printed p.~20, says that the braid diagram has a
decomposition circle if and only if some seesaw number belongs to
$\{2,3\}$. By Lemma~\ref{lem:seesaws} and $n\ge2$, all the seesaw
numbers here are at least four. There is no decomposition circle, hence
the knot is prime. Its positive genus below also excludes the unknot.
\end{proof}

\begin{proof}[Genus and crossing assertions of Theorem~\ref{thm:family}]
The word $W_b^n$ has $n[17+3(b-4)]=n(3b+5)$ crossings. Its oriented
smoothing has $b$ Seifert circles. Its canonical Seifert surface is
connected, since every adjacent braid generator occurs. Its genus is
therefore
\[
                  \frac{n(3b+5)-b+1}{2}.
\]
For a homogeneous braid closing to a knot, its diagram is homogeneous
\cite[Lemma~2.2, printed p.~5]{abe}. Cromwell's minimal-genus theorem
for homogeneous diagrams, stated in \cite[Theorem~2.3, same page]{abe},
says that this canonical surface realizes the knot genus.

Yamada's theorem \cite[Theorem~3, printed p.~347]{yamada} identifies the
minimum number of Seifert circles over all diagrams with the braid
index, already proved to be $b$. Any diagram $D$ of this knot has
$g_D\ge g(K_{b,n})$ and at least $b$ Seifert circles. The Seifert
algorithm identity consequently gives
\[
 c(D)=2g_D+s(D)-1\ge2g(K_{b,n})+b-1=n(3b+5).
\]
The displayed diagram attains equality. For fixed $b$, take
$n=kb+1$ with $k\ge1$; the resulting crossing numbers strictly increase,
so these are infinitely many distinct knots.
\end{proof}

\begin{corollary}\label{cor:four}
For every odd positive $n$, the closure of $\beta^n$ is prime, with
\[
 \mr=\br=\bi=4,\qquad g=\frac{17n-3}{2},\qquad c=17n.
\]
It has the same meridian-class epimorphism onto $A_8$.
\end{corollary}
\begin{proof}
The braid permutation is a four-cycle, so odd $n$ gives a knot. The
group and genus arguments above apply unchanged. The only seesaw
numbers are $4n$ and $6n+1$, both at least four even when $n=1$.
The same visual-primeness argument applies.
\end{proof}

\section{A connected-sum corollary with simple targets}

Let $K_*=\widehat\beta$ and let $3_1$ denote a trefoil of the chirality
represented by the positive two-braid $\sigma_1^3$.

\begin{corollary}\label{cor:sums}
Let $K$ be any knot with a bridge presentation using at most four bridges
and an epimorphism onto $A_8$ taking meridians to double transpositions.
For every $s\ge0$, the knot $J_s=K\mathbin{\#}(\mathbin{\#}^{s}3_1)$
has a meridian-class epimorphism onto $A_{8+2s}$ and
\[
                        \mr(J_s)=\br(J_s)=4+s.
\]
The target has ordinary generating rank two and double-transposition
class generating rank $4+s$. In particular, this holds for $K=K_*$.
\end{corollary}
\begin{proof}
Embed the given target as $A_8$ on the first eight points. Since its
double transpositions form one conjugacy class, conjugate the
homomorphism so that a selected meridian maps to $x_3$.
A trefoil group has a two-meridian presentation
$\langle a,b\mid aba=bab\rangle$. For two involutions $u,v$, the
condition $(uv)^3=1$ implies $uvu=vuv$. Glue the meridian $a$ to the
last chosen meridian of the preceding summand and map
$(a,b)$ to $(x_{j-1},x_j)$, for $j=4,\ldots,3+s$.
The connected-sum group is the amalgam over these meridians; the
compatible homomorphisms therefore define a homomorphism to $A_{8+2s}$.
The images generate that whole group by Theorem~\ref{thm:targets}.
The class-rank bound gives $\mr(J_s)\ge4+s$. Joining bridge positions
with at most four bridges for $K$ and two for each trefoil removes one bridge
per join, giving $\br(J_s)\le4+s$. The elementary inequality
$\mr\le\br$ proves equality. No statement about branched covers is a
premise of this argument.
\end{proof}

The same braid gluing identifies $\widehat{W_b}$ with
$K_*\mathbin{\#}(\mathbin{\#}^{b-4}3_1)$. Thus the power restriction
$n\ge2$ in Theorem~\ref{thm:family} has content: at $n=1$ the later
seesaws are two and the diagrams display these connected sums.

\section{Literal control, replay and scope}\label{sec:literal}

The supplement
\path{alternating_target_prime_families_inputs.json} contains the signed
literal braid, all forty PD rows and all forty arc permutations for
$K_{5,2}$. It retains the five seed arcs and thirty-five successive
Wirtinger coloring moves. Arcs are the components obtained by joining
the two overpassing ports in each PD row, ordered lexicographically by
their edge-label sets. For a PD row $(a,b,c,d)$ the overarc contains
$b,d$, and the underarcs contain $a,c$. The signed relation is
\[
 p_c=p_b^{\epsilon}p_a p_b^{-\epsilon},\qquad\epsilon\in\{1,-1\};
\]
all recorded images are involutions, so both signs are checked by the
same exact permutation identity. The signs and braid port wiring are
retained separately to bind this certificate to the literal oriented
braid, rather than only to an abstract tuple.

The standard-library checker
\path{verify_alternating_target_prime_families.py} verifies the two
short base identities, every base braid step, the induction identities,
the full signed PD and closure wiring, every Wirtinger relation, the
five-seed trace, and exact adjacent-generator seesaws in finite controls.
It checks the spherical rotation system and the one-component condition
of the literal diagram. An optional topology replay reconstructs the
forty-row PD directly from the five-braid using Spherogram. Corrupted
base words, arc colors, PD wiring and extension generators are rejected,
as are invalid theorem parameters. The replay hashes the three primary
source PDFs cited below. These finite controls supplement the written
proof of the uniform assertions; they are not a search over all $b,n$.

There is a useful positive comparison control at $(b,n)=(5,2)$.
On six points the transposition tuple
\[
 ((0\,1),(0\,2),(0\,3),(2\,4),(0\,5))
\]
also returns under $W_5^2$. Its edges form a connected tree, so it
generates $S_6$. The checker screens all $216$ simultaneous-conjugacy
orbits of transitive five-transposition tuples on six points; five
orbits return. This establishes a symmetric quotient for this member,
and supplies no absence claim for arbitrary Coxeter targets. Neither
augmentation-rank values nor exclusions from every previously studied
knot family are conclusions of this paper.

\paragraph{AI assistance and review.}
ChatGPT (OpenAI) assisted with mathematical arguments, literature checks,
exposition, proof review and verification code.

\begin{thebibliography}{9}
\bibitem{flo}
P.~Feller, L.~Lewark and M.~Orbegozo Rodriguez,
\emph{Homogeneous braids are visually prime},
Journal of Topology \textbf{18} (2025), e70040,
\href{https://doi.org/10.1112/topo.70040}{doi:10.1112/topo.70040};
\href{https://arxiv.org/abs/2408.15730}{arXiv:2408.15730v2},
30 September 2025, Theorem~1.1 and Lemma~5.1(i).
\bibitem{abe}
T.~Abe, \emph{The Rasmussen invariant of a homogeneous knot},
Proceedings of the American Mathematical Society \textbf{139} (2011),
2647--2656; \href{https://arxiv.org/abs/1003.5392}{arXiv:1003.5392v1},
Lemma~2.2 and Theorem~2.3 (Cromwell's minimal-genus theorem).
\bibitem{yamada}
S.~Yamada, \emph{The minimal number of Seifert circles equals the braid
index of a link}, Inventiones Mathematicae \textbf{89} (1987), 347--356,
\href{https://doi.org/10.1007/BF01389082}{doi:10.1007/BF01389082},
Theorem~3.
\end{thebibliography}
\end{document}
